\section{BREW-RC}
\label{sec:brewrc-method} Section~\ref{sec:brewrc-exp4} demonstrates
that write latency distributions identify the polarity category of
entire codeword transitions. These distributions are labeled as:
\begin{equation}
	\ell_{cw},\, \ell_d \in \{\texttt{NONE},\, \texttt{UNIPOLAR\_P},\, \texttt{UNIPOLAR\_N},\, \texttt{BIPOLAR}\},
\end{equation}
where $\ell_{cw}$ reflects the codeword DHD and $\ell_d$ is the known transition polarity of the applied data. The relationship between the applied data label and the observed codeword label is written as $\ell_d \to \ell_{cw}$.
Four observation types arise. $\texttt{NONE} \to \texttt{NONE}$ transitions reveal nothing about the parity function; as shown in Section~\ref{sec:brewrc-exp1}, their timing is governed by write-path serialization overhead with near-zero variance. $\texttt{BIPOLAR} \to \texttt{BIPOLAR}$ observations are similarly uninformative, as a forced bipolar data transition permits parity bits to transition in any polarity. Observations of $\texttt{NONE} \to \texttt{UNIPOLAR}_{N/P}$ or $\texttt{BIPOLAR} \to \texttt{UNIPOLAR}_{N/P}$ would violate the hypothesis that parity bit activity is the source of timing fluctuations; assertions confirm these are never observed.
BREW-RC uses two observation classes to generate SAT constraints. $\texttt{UNIPOLAR}_{N/P} \to \texttt{UNIPOLAR}_{N/P}$ observations produce explicit constraints: given a unipolar data transition, no parity bits transitioned in an opposing polarity. $\texttt{UNIPOLAR}_{N/P} \to \texttt{BIPOLAR}$ observations produce existential constraints: given a unipolar data transition, at least one parity bit transitioned in an opposing polarity.

\subsection{Modeling Parity-Bit Activity}
For a dataword $d \in \{0,1\}^k$, the codeword $c = dG = [d \mid p]$
appends parity bits generated by the $\mathrm{GF}(2)$ multiplication:
\begin{equation}
	p = dP.
\end{equation}
For pre- and post-write states $(d_0, d_1)$, the parity-bit updates are:
\begin{equation}
	\Delta p = (d_0 + d_1)P = \Delta d\, P,
\end{equation}
where $\Delta d = d_0 + d_1$ marks toggled data bits. $\Delta p$ acts as a gate that determines which parity bits update under a given data transition, independent of transition polarity. In the SAT construction, $\Delta p$ is evaluated for every dataword transition and used in all clauses to restrict parity-bit activity to the toggled subset.

\subsection{Polarity-Encoded SAT Formulation}
\label{sec:sat-encoding} Each transition is assigned a data-polarity
label $\ell_d \in \{\texttt{UNIPOLAR\_P}, \texttt{UNIPOLAR\_N}\}$ with
scalar representation $d \in \{-1, +1\}$: $d=+1$ for positive (${\to}1$)
and $d=-1$ for negative ($1{\to}0$) unipolar transitions. Bipolar and
static cases (\texttt{BIPOLAR}, \texttt{NONE}) produce no directional
information and are excluded. For each parity bit $p_j$, Boolean
variables $p_{0,j}$ and $p_{1,j}$ denote its pre- and post-write values,
gated by $\Delta p_j$. Positive and negative parity transitions are
encoded as:
\begin{align}
	\textit{pos}_j & \leftrightarrow \Delta p_j \wedge (\neg p_{0,j} \wedge p_{1,j}), \\
	\textit{neg}_j & \leftrightarrow \Delta p_j \wedge (p_{0,j} \wedge \neg p_{1,j}),
\end{align}
which encode positive and negative parity transitions only when the parity bit is expected to update. Depending on $(\ell_d, \ell_{cw})$, one of four clause types is generated as summarized in Table~\ref{tab:polarity-clauses}.
\begin{table}[t!]
\setlength{\textfloatsep}{0pt}
\setlength{\floatsep}{0pt}
\centering
\caption{SAT solver clause definitions \label{tab:polarity-clauses}.}
\vspace{-.2cm}
\renewcommand{\arraystretch}{1.1}
\setlength{\tabcolsep}{3pt}
\begin{tabular*}{\linewidth}{@{\extracolsep{\fill}}lcccc@{}}
\toprule
\textbf{Data label} $\ell_d$ & \textbf{Obs.} $\ell_{cw}$ & \textbf{Constraint} & \textbf{Description} \\
\midrule
\texttt{UNIPOLAR\_P} & \texttt{UNIPOLAR\_P} & $\forall j\, (p_{0,j}\vee\neg p_{1,j})$  & forbid \emph{neg}  \\
\texttt{UNIPOLAR\_N} & \texttt{UNIPOLAR\_N} & $\forall j\, (\neg p_{0,j}\vee p_{1,j})$ & forbid \emph{pos} \\
\texttt{UNIPOLAR\_P} & \texttt{BIPOLAR}     & $\bigvee_j \textit{neg}_j$               & require $\ge1$ \emph{neg}  \\
\texttt{UNIPOLAR\_N} & \texttt{BIPOLAR}     & $\bigvee_j \textit{pos}_j$               & require $\ge1$ \emph{pos}  \\
\bottomrule
\end{tabular*}  
\end{table}
Each $(d_0, d_1, \ell_d, \ell_{cw})$ tuple contributes a compact SAT subproblem constraining the allowed polarity of parity-bit transitions under a candidate parity submatrix $P$.

\subsection{Unipolar Random Word Walk}
Because the preliminary experiments exercised only single-bit and
single-byte toggles, the resulting clauses may not reveal cross-byte
parity coupling between data bits in different bytes. A \emph{unipolar
	random word walk} is introduced to systematically explore unipolar
transitions over full word regions, efficiently sampling the word-level
unipolar transition space. The procedure is described in
Algorithm~\ref{alg:ruw}.
\setlength{\textfloatsep}{-10pt} {
	% \setlength{\floatsep}{-10pt}
	\begin{algorithm}[t!]
		\footnotesize
		\caption{1-word, random, unipolar walk\label{alg:ruw}}
		\begin{algorithmic}[1]
			\Require word size $k$, length $\ell$, max test-vectors per address $\textit{max\_tv}$
			\State $\texttt{Keys} \gets \varnothing$; $\texttt{A} \gets \varnothing$; \Comment{used transitions, used addrs.}
			\While{$i < \ell$}
			\If{$i \bmod \text{max\_tv} = 0$}
			\State $\texttt{addr}\gets\texttt{GetRandWordAddr(A)}$; $\texttt{A} \cup \texttt{addr}$;
			\State $s \gets \texttt{GetRandWord()}$;
			\State write\_word(addr, $s$)
			\EndIf
			\State $d \gets  s=0 \; ? \; 1 : s=2^k-1 \; ? \; 0 : \texttt{RNG(0,1)}))$;
			\State $m \gets \texttt{RNG(1, $2^k -1$)}$
			\State $s' \gets d=0\; ? \;s\; \& \; \neg m : s\; | \; m $;
			\State $key \gets \min(s,s') << k \; | \; max(s, s')$;
			\If{$\texttt{key} \notin \texttt{Keys}$}
			\State $\texttt{Keys} \cup \texttt{k}$ ; $s \gets s'$ ; $i\gets i+1$;
			\State write\_word($\texttt{addr}$, $s'$)
			\EndIf
			\EndWhile
		\end{algorithmic}
	\end{algorithm}
}
\setlength{\textfloatsep}{3pt}

\subsection{Data Collection and SAT Solving}
\label{sec:brewrc-datacollection} To ensure every test vector emits a
useful SAT constraint, a random unipolar walk of the dataword space is
performed and $t_w$ is measured for every applied test vector. The 4M
has a higher probability of generating $\texttt{UNIPOLAR}_{P/N} \to
	\texttt{UNIPOLAR}_{P/N}$ constraints due to its smaller parity
submatrix. For each chip, the experiment is repeated 10 times at
different word addresses to generate a representative $t_{wmin}$. To
converge to a single ECC solution, 100K (200K) $t_{wmin}$ measurements
are collected from each 4M (8M) sample. After capturing timing data,
labels are assigned to construct $(d_0, d_1, \ell_d, \ell_{cw})$
observations. Each observation yields a set of XOR constraints encoding
$\Delta p$ and CNF clauses enforcing the allowed polarity of each
parity-bit transition. The resulting hybrid XOR+CNF formula is solved
using CryptoMiniSAT~\cite{soos2009extending, soos2018cryptominisat5}.
Since BREW-RC cannot infer the number of parity bits $r$ a priori, the
solver is run for $r \in \{1, \ldots, k\}$, corresponding to all
candidate matrices up to a code rate of 0.5. BREW-RC returns a single
parity submatrix satisfying a unique $r$ for both the 4M and 8M.
